\chapter{椭圆曲线与模形式的 L 函数}\label{ch:05}

\begin{center}
Ehud de Shalit
\end{center}

上一章考察了整体域的 Galois 群到 $\GL(n,\mathbb{C})$ 的表示及其 L 函数。这类表示必然具有有限像。我们在那一章结尾看到，权 $1$ 的模形式提供了重要的例子。

然而，算术中出现的 $G_{\mathbb{Q}}=\operatorname{Gal}(\overline{\mathbb{Q}}/\mathbb{Q})$ 的表示大多没有有限像，即不能通过 $\mathbb{Q}$ 的有限扩张分解。粗略地说，取一有限组系数在 $\mathbb{Q}$ 中的多项式方程，设它们定义维数为 $d$ 的代数簇；则对每个指标 $0\le i\le 2d$，étale 上同调理论都会附上一族 $\ell$ 进表示（$G_{\mathbb{Q}}$ 到 $\GL_n(\mathbb{Q}_\ell)$ 的表示，$\mathbb{Q}_\ell$ 为 $\ell$ 进数域），每个 $\ell$ 一个。Artin 表示对应 $d=0$。它们具有有限像，本质上与 $\ell$ 无关，因而可在任何特征 $0$ 的代数闭域中实现，特别地可在 $\mathbb{C}$ 中实现。把 $\mathbb{Q}$ 换成任意整体域，结论同样成立。

与上一章一样，可以对每个这样的 $\ell$ 进表示附上一个 L 函数。它是由 Euler 乘积给出的 Dirichlet 级数，在某个右半平面内收敛，系数应当是代数的，依赖于指标 $i$ 而不依赖 $\ell$。在数域情形，这些 L 函数还要补上阿基米德位处的 Euler 因子。它们应当满足解析延拓与函数方程，形式与 Artin L 函数的相仿。一般而言，这些都是宏大的未决猜想。

这两章讨论 Artin L 函数之外第一个有趣的情形。我们把基域限制为 $\mathbb{Q}$（放弃函数域这个非常重要的例子），并取代数簇为一条曲线（$d=1$），而且取仍能产生非平凡结果的最简单曲线：椭圆曲线（亏格 $1$ 的曲线）。指标 $i$ 取 $1$（$i=0$ 或 $2$ 给出 $\zeta(s)$ 与 $\zeta(s+1)$；一般只有范围 $0<i<2d$ 有趣，而在很多情形——例如射影空间中的超曲面——唯一有趣的指标是 $i=d$）。

\begin{annote}
原文此处写“$i=0$ 或 $2$ 给出 $\zeta(s)$ 与 $\zeta(s+1)$”，其中 $\zeta(s+1)$ 应为 $\zeta(s-1)$。对 $\mathbb{Q}$ 上光滑射影曲线 $C$，$\ell$ 进上同调 $H^2(C_{\overline{\mathbb{Q}}},\mathbb{Q}_\ell)$ 是一维的 Tate 扭转 $\mathbb{Q}_\ell(-1)$，其 L 函数为 $\zeta(s-1)$；$H^0$ 才给出 $\zeta(s)$。$H^1$ 贡献曲线自身的 L 函数，正是本章 $i=1$ 的情形。原书笔误照录，读者勿据以计算。
\par\vspace{0.15em}{\raggedleft 【来源：原书 p.~89（A）；Milne《Lectures on Étale Cohomology》（A）；Silverman《The Arithmetic of Elliptic Curves》附录 C（A）】\par}
\end{annote}

定义在 $\mathbb{Q}$ 上的椭圆曲线所附的 $\ell$ 进表示是奇的、二维的、不可约的。与 Artin L 函数的情形一样，这类表示的重要来源由模形式提供，此处是权 $2$ 的模形式；Weil 的逆定理预言它们全都来自模形式。这就是著名的 Shimura--Taniyama--Weil 猜想，它也可用椭圆曲线的模一致化来表述。它现在已是定理，归功于 A.~Wiles 与 Taylor--Wiles 的工作，以及 F.~Diamond、R.~Taylor、B.~Conrad、C.~Breuil 后来的补充工作。

\section{椭圆曲线}\label{c5:sec:1}

\subsection{$\mathbb{C}$ 上的椭圆曲线}\label{c5:sec:1.1}

复椭圆曲线是亏格 $1$ 的紧 Riemann 曲面 $E$。由单值化定理，它同构于
\begin{equation}\tag{1.1}\label{c5:eq:1.1}
E_\tau=\mathbb{C}/(\mathbb{Z}+\mathbb{Z}\tau),
\end{equation}
其中 $\tau\in\mathfrak{H}$ 是复上半平面中的点，且 $\tau$ 在 $\mathrm{SL}(2,\mathbb{Z})$ 的作用下唯一，作用为
\begin{equation}\tag{1.2}\label{c5:eq:1.2}
\begin{pmatrix}a&b\\ c&d\end{pmatrix}:\tau\longmapsto\frac{a\tau+b}{c\tau+d}.
\end{equation}

$E$ 上亚纯函数构成的域（即椭圆函数域）$\mathbb{C}(E)$ 的超越次数为 $1$。Weierstrass 写出了它的生成元。设 $\mathcal{L}=\mathbb{Z}+\mathbb{Z}\tau$，或者更一般地，设 $\mathcal{L}$ 为 $\mathbb{C}$ 中任意格（它当然与上述形状的某个格位似），并定义
\begin{equation}\tag{1.3}\label{c5:eq:1.3}
\wp(z,\mathcal{L})=\frac{1}{z^2}+\sum_{\omega\in\mathcal{L}-\{0\}}\Bigl(\frac{1}{(z-\omega)^2}-\frac{1}{\omega^2}\Bigr).
\end{equation}

\begin{annote}
$\wp$ 的定义中减去 $1/\omega^2$ 并非装饰。若只写 $\sum 1/(z-\omega)^2$，则当 $|\omega|\to\infty$ 时通项为 $2z/\omega^3+O(|\omega|^{-4})$，级数并不绝对收敛。减去 $1/\omega^2$ 后通项是 $O(|\omega|^{-3})$，在 $\mathbb{C}\setminus\mathcal{L}$ 的紧集上一致收敛。这个正规化也保证了 $\wp$ 的 $\mathcal{L}$ 周期性：逐项验证 $\wp(z+\omega_0)-\wp(z)$ 的导数为零，再用偶函数性定出该差为零。另一常见约定是把和写成 $\sum_{\omega\neq0}\bigl((z-\omega)^{-2}-\omega^{-2}\bigr)$，与 (1.3) 完全一致。
\par\vspace{0.15em}{\raggedleft 【来源：Silverman《The Arithmetic of Elliptic Curves》第 VI 章（A）；Lang《Elliptic Functions》第 1 章（A）】\par}
\end{annote}

于是 $\wp$ 收敛、亚纯、以 $\mathcal{L}$ 为周期，它在 $E$ 上唯一的极点是 $0\pmod{\mathcal{L}}$ 处的二阶极点；$\wp$ 与其导数 $\wp'$ 一起生成 $\mathbb{C}(E)=\mathbb{C}(\wp,\wp')$。由于这是超越次数 $1$ 的域，$\wp$ 与 $\wp'$ 代数相关。事实上它们满足 Weierstrass 方程，该方程可看作 $\wp$ 的微分方程：
\begin{equation}\tag{1.4}\label{c5:eq:1.4}
\wp'^2=4\wp^3-g_2(\mathcal{L})\wp-g_3(\mathcal{L}),
\end{equation}
其中常数由权 $4$ 与权 $6$ 的 Eisenstein 级数给出：
\begin{equation}\tag{1.5}\label{c5:eq:1.5}
g_2(\mathcal{L})=60\sum_{\omega\in\mathcal{L}-\{0\}}\omega^{-4},\qquad
g_3(\mathcal{L})=140\sum_{\omega\in\mathcal{L}-\{0\}}\omega^{-6}.
\end{equation}

作为抽象域，
\begin{equation}\tag{1.6}\label{c5:eq:1.6}
\mathbb{C}(E)\cong\mathbb{C}(X)[Y]/(Y^2-4X^3+g_2X+g_3).
\end{equation}
$X$ 中三次方程的判别式为
\begin{equation}\tag{1.7}\label{c5:eq:1.7}
\Delta(\mathcal{L})=g_2^3-27g_3^2\neq0.
\end{equation}

\begin{annote}
(1.7) 把 $\Delta$ 定义为 $g_2^3-27g_3^2$，而 (1.6) 中的三次式是 $Y^2=4X^3-g_2X-g_3$。按三次多项式判别式的标准公式，$4X^3-g_2X-g_3$ 的判别式是 $16(g_2^3-27g_3^2)$，故 (1.7) 比它小一个因子 $16$。另一方面，若把 (1.4) 写成标准形式 $y^2=x^3+ax+b$，则 $a=-g_2/4$、$b=-g_3/4$，其 Weierstrass 判别式恰为 $g_2^3-27g_3^2$，与 (1.7) 一致。可见因子 $16$ 的差别来自“三次式的判别式”与“Weierstrass 方程的判别式”两个不同约定；本章只用 $\Delta\neq0$，两种约定等价。
\par\vspace{0.15em}{\raggedleft 【来源：Silverman《The Arithmetic of Elliptic Curves》第 III 章（A）；原书 (1.4)--(1.7)（A）】\par}
\end{annote}

反之，每个判别式非零的 Weierstrass 方程都来自某个格 $\mathcal{L}$。（警告：若两个格位似，不同的方程可能表示同构的椭圆曲线。）

嵌入 $E-\{0\}\hookrightarrow\mathbb{C}^2\subset\mathbb{P}^2(\mathbb{C})$ 由
\begin{equation}\tag{1.8}\label{c5:eq:1.8}
z\longmapsto(\wp(z),\wp'(z))=(\wp(z):\wp'(z):1)
\end{equation}
给出，它可延拓为 $E$ 作为 $\mathbb{P}^2(\mathbb{C})$ 中光滑射影曲线的嵌入。唯一的“无穷远点” $(0:1:0)$ 是 $0_E=0\bmod\mathcal{L}$ 的像。

把 $E$ 描述为复环面，立刻看出它带有以 $0_E$ 为原点的阿贝尔群结构。用射影嵌入来计算群律是有教益的。设 $P$ 与 $P'$ 是两点，代数几何的初等考虑给出群律的下述解释：设 $Q$ 是过 $P$ 与 $P'$ 的直线（若 $P=P'$ 则取 $P$ 处的切线）与三次曲线 $E$ 的第三个交点（回忆 Bézout 定理！）。则 $P+P'$ 是把 $Q$ 关于 $x$ 轴反射得到的点。

\subsection{任意域上的椭圆曲线}\label{c5:sec:1.2}

在任意域 $F$ 上，椭圆曲线是亏格 $1$ 的光滑、射影、几何不可约曲线 $E$，并带有一个 $F$ 有理点 $0_E$。由 Riemann--Roch 定理，$E$ 带有阿贝尔群簇的结构，运算定义在 $F$ 上，$0_E$ 是单位元。

若 $F$ 是 $\mathbb{C}$ 的子域，则复椭圆曲线可在 $F$ 上定义，只要其 Weierstrass 方程可选得系数在 $F$ 中。但要注意：$F$ 上不同构的椭圆曲线可能在 $\mathbb{C}$ 上同构！此时群运算的 $F$ 有理性由上述三次曲线的几何描述立刻看出。更一般地，特征不等于 $2$ 或 $3$ 的域上的椭圆曲线有平面射影模型，由 $g_2,g_3\in F$ 的 Weierstrass 方程给出。Tate 也把“Weierstrass 模型”推广到了特征 $2$ 与 $3$。

对 $F$ 的任意扩域 $K$，$E$ 的 $K$ 有理点记作 $E(K)$，构成阿贝尔群。若 $K$ 有限，它是有限群。若 $K$ 是整体域（更一般地，是其素域上有限生成的域），它是有限生成阿贝尔群。（这相当深刻：即 Mordell--Weil 定理。）其（自由）秩称为 $E(K)$ 的 Mordell--Weil 秩。

椭圆曲线典范同构于其雅可比簇，也同构于其 Picard 簇（分类次数 $0$ 线丛的簇，或同样地，次数 $0$ 除子模主除子）。到 Picard 簇的同构由把 $x\in E(F)$ 送到 $(x)-(0)$ 的类（模主除子）实现。Abel--Jacobi 定理断言：除子 $D=\sum n_i(x_i)$ 是主的，当且仅当 $\deg(D)=\sum n_i=0$ 且 $s(D)=\sum n_ix_i=0_E$。在 $\mathbb{C}$ 上可用解析方法证明，即 Cauchy 积分公式与 theta 函数。代数证明则由 Riemann--Roch 定理推出。

\subsection{自同态环}\label{c5:sec:1.3}

椭圆曲线的同源是非零的群簇同态 $\alpha\colon E\to E'$。其次数 $\deg(\alpha)$ 定义为 $E'$ 的一般点上方纤维中点的个数（按重数计，若 $\alpha$ 不可分），等价地
\begin{equation}\tag{1.9}\label{c5:eq:1.9}
\deg(\alpha)=[F(E):\alpha^*F(E')].
\end{equation}
次数是正整数。把 $\operatorname{Ker}(\alpha)$ 记作 $E[\alpha]$。它是有限群概形，既约当且仅当 $\alpha$ 可分。在纯不可分情形，它只是“被幂零元加厚”的原点。$\alpha$ 的次数等于它作为群概形的阶。在可分情形，它就是 $E[\alpha]$ 在 $\overline{F}$ 中点的实际个数。

若 $\alpha\colon E\to E'$ 是同源，对偶同源 $\check{\alpha}\colon E'\to E$ 是由 Picard 函子性得到的同源，用到椭圆曲线是其自身的 Picard 簇。它也由下述关系唯一刻画：
\begin{equation}\tag{1.10}\label{c5:eq:1.10}
\check{\alpha}\circ\alpha=\deg(\alpha)\cdot 1_E.
\end{equation}

$E$ 在 $F$ 上的自同态环 $\operatorname{End}(E)$ 按定义是 $E$ 的全部代数群自同态构成的环。（它们构成环，因为 $E$ 交换。）一般它不会比 $\mathbb{Z}$ 大。若更大，称 $E$ 具有复乘（CM）。若 $\operatorname{char}(F)=0$，于是可设 $F\subset\mathbb{C}$，并像前面那样给 $E$ 附上一点 $\tau\in\mathfrak{H}$（模 $\mathrm{SL}(2,\mathbb{Z})$ 的作用），则 $\operatorname{End}(E)$ 或者是 $\mathbb{Z}$，或者是二次虚域 $\mathbb{Q}(\tau)$ 中的一个序，只要后者确实是二次虚域。由复单值化立刻看出：$\operatorname{End}(E/\mathbb{C})$ 就是满足 $\alpha\mathcal{L}\subset\mathcal{L}$ 的复数 $\alpha$ 的集合。这也解释了“复乘”一词的来历！

若 $\operatorname{char}(F)=p>0$，自同态环可能更大，且非交换。此时它是 $\mathbb{Q}$ 上唯一四元数代数中的极大序，该四元数代数在 $p$ 与 $\infty$ 处分歧，在别处分裂。这样的椭圆曲线称为超奇异的，它们只有有限多条。否则 $E$ 称为常的。无论如何，作为阿贝尔群，自同态环是秩 $1$、$2$ 或 $4$ 的自由群。这些结果属于 Deuring，可追溯到 1930 年代。证明要借助 $E$ 的 $\ell$ 进 Tate 模（见下节）。

可以证明 $\deg$ 是（自由）阿贝尔群 $\operatorname{End}(E)$ 上的二次函数。这需要一些工夫，尤其当 $\operatorname{char}(F)>0$；正确的做法是用立方定理（见 Mumford 的书《Abelian varieties》，其中讨论了阿贝尔簇自同态环的更一般情形）。于是 $\deg(n\cdot1-\alpha)$ 是 $n$ 的首一二次多项式，系数为整数，我们把它写成
\begin{equation}\tag{1.11}\label{c5:eq:1.11}
n^2-\operatorname{tr}(\alpha)n+\deg(\alpha).
\end{equation}
它称为 $\alpha$ 的特征多项式。例如，若 $F$ 特征为 $0$，且 $E$ 具有序 $\operatorname{End}(E)\subset\mathbb{Q}(\tau)$ 的复乘，则次数与 $\mathbb{Q}(\tau)$ 到 $\mathbb{Q}$ 的范数一致。

一个非常重要的情形是 $F=\mathbb{F}_q$ 为有限域。例如，用 Weierstrass 方程来想，看出映射 $\varphi\colon E\to E$，
\begin{equation}\tag{1.12}\label{c5:eq:1.12}
\varphi(x,y)=(x^q,y^q)
\end{equation}
是曲线的自同态（注意：虽然它在 $E(\overline{\mathbb{F}}_q)$ 上是双射，却不是自同构；它是次数为 $q$ 的纯不可分自同态）。这个特殊的自同态称为 $E$ 的 Frobenius 自同态（次数为 $q$）。

$\operatorname{End}(E)$ 与 $\operatorname{Gal}(\overline{F}/F)$ 都作用在 $E(\overline{F})$ 上，且两个作用交换。一般地它们非常不同；例如，若 $F$ 是数域，自同态连续作用（在复拓扑下），但除复共轭以外的 Galois 元作用不连续。

然而，若 $F=\mathbb{F}_q$ 有限，$\varphi$ 的几何作用与 $\sigma$ 的相同，这里 $\sigma$ 是 $\operatorname{Gal}(\overline{\mathbb{F}}_q/\mathbb{F}_q)$ 的生成元，把每个域元素提升到 $q$ 次幂。由于由 Galois 理论，$\sigma^n$ 的不动点是 $E(\mathbb{F}_{q^n})$ 的点，我们得到
\begin{equation}\tag{1.13}\label{c5:eq:1.13}
E(\mathbb{F}_{q^n})=\operatorname{Fix}(\varphi^n)=E(\overline{F})[\varphi^n-1].
\end{equation}
自同态 $\varphi^n-1$ 可分，故其次数等于其核中元素的个数。由次数的二次公式得到：

\par\medskip\noindent\textbf{推论 1.1}\quad 设 $F=\mathbb{F}_q$。则
\begin{equation}\tag{1.14}\label{c5:eq:1.14}
\#E(\mathbb{F}_{q^n})=1-(\lambda^n+\lambda'^n)+q^n,
\end{equation}
其中 $\lambda$ 与 $\lambda'$ 是 $\varphi$ 的特征多项式的两个根。

这是因为 $\varphi^n$ 的特征多项式的根为 $\lambda^n$ 与 $\lambda'^n$（一旦建立了特征多项式与 $\ell$ 进表示的联系，这一点便随之成立；见下文），且其次数为 $q^n$。

一个非常重要的观察是：二次型 $\deg(n-m\varphi)$ 正定，除非 $\varphi$ 是整数，此时它半正定。在后一情形 $\lambda=\lambda'=\sqrt{q}$ 是整数。否则 $\lambda$ 与 $\lambda'$ 互为复共轭，生成一个二次虚域。特别地 $|\lambda|=|\lambda'|=\sqrt{q}$。于是有下述定理。

\par\medskip\noindent\textbf{定理 1.2}\quad
\begin{equation}\tag{1.15}\label{c5:eq:1.15}
|1+q-\#E(\mathbb{F}_q)|\le2\sqrt{q}.
\end{equation}

这就是 Hasse 的博士论文（1934），由 E.~Artin 猜想。我们将会看到，它等价于 $E/\mathbb{F}_q$ 的 Riemann 假设。

\section{$\ell$ 进表示}\label{c5:sec:2}

\subsection{Tate 模}\label{c5:sec:2.1}

在任意域 $F$ 上，若 $\ell$ 与 $\operatorname{char}(F)$ 互素，乘以 $\ell$ 是可分自同态，故 $E[\ell]$ 是阶为 $\ell^2$、指数为 $\ell$ 的阿贝尔群。因此它同构于 $(\mathbb{Z}/\ell\mathbb{Z})^2$。同样可证
\begin{equation}\tag{2.1}\label{c5:eq:2.1}
T_\ell E=\varprojlim E[\ell^n]\cong\mathbb{Z}_\ell^2.
\end{equation}
这个模称为 $E$ 的 Tate 模。若 $p=\operatorname{char}(F)>0$，对 $T_pE$ 用同样的定义，不难看出此时若 $E$ 是常的，则 $T_pE\cong\mathbb{Z}_p$，若 $E$ 超奇异则为 $0$。

\begin{annote}
特征 $p$ 中 $T_pE$ 的行为与 $\ell\neq p$ 时完全不同。$E[p]$ 是阶为 $p$ 的连通群概形，其 $p$ 可除点系统在逆向极限中消失。若 $E$ 超奇异，$E[p]$ 连 $p$ 阶 étale 部分都没有，故 $T_pE=0$；若 $E$ 常，$E[p]$ 有 étale 商 $\mathbb{Z}/p$，故 $T_pE\cong\mathbb{Z}_p$。因此不能把“$T_\ell E$ 是秩 $2$ 自由 $\mathbb{Z}_\ell$ 模”这句话照搬到 $\ell=p$。本章构造 L 函数只用 $\ell\neq p$ 的 $T_\ell E$，$T_pE$ 只在讨论好归约判别法时作为反例出现。
\par\vspace{0.15em}{\raggedleft 【来源：Silverman《The Arithmetic of Elliptic Curves》第 III、V 章（A）；Tate《$p$-divisible groups》（A）】\par}
\end{annote}

$\operatorname{End}(E)$ 与 $G_F=\operatorname{Gal}(\overline{F}/F)$ 都连续作用在 $T_\ell E$ 上，且两个作用交换。顺便说，Deuring 正是这样证明了他关于 $\operatorname{End}(E)$ 的结构定理。注意自同态环的作用是忠实的，即使在 $T_pE$ 上（只要它不为 $0$），因为非零自同态不可能有无穷大的核。

关于 $\operatorname{End}(E)$ 的作用，有下述定理。

\par\medskip\noindent\textbf{定理 2.1}\quad 设 $\alpha\in\operatorname{End}(E)$，$T_\ell(\alpha)$ 是它在 $T_\ell E$ 上作用的矩阵。则
\begin{equation}\tag{2.2}\label{c5:eq:2.2}
\det T_\ell(\alpha)=\deg(\alpha).
\end{equation}

\par\medskip\noindent\textbf{推论 2.2}\quad $T_\ell(\alpha)$ 的特征多项式等于 $\alpha$ 的特征多项式。特别地，它系数为整数，且与 $\ell$ 无关。

\begin{annote}
(2.2) 只给出行列式；要得到整系数的特征多项式，还需迹的整数性。由 (1.11) 取 $n=1$ 得 $\operatorname{tr}(\alpha)=1+\deg(\alpha)-\deg(1-\alpha)$，右端全是整数，故 $\operatorname{tr}T_\ell(\alpha)$ 也是整数。这正是推论 2.2 的实质：$\alpha$ 的特征多项式 $X^2-\operatorname{tr}(\alpha)X+\deg(\alpha)$ 与 $\ell$ 无关。若只证行列式，$T_\ell(\alpha)$ 的特征多项式仍可能随 $\ell$ 变动；迹的整数性把这一点堵死。
\par\vspace{0.15em}{\raggedleft 【来源：Silverman《The Arithmetic of Elliptic Curves》第 III 章（A）；原书 (1.11) 与定理 2.1（A）】\par}
\end{annote}

习惯上把 Galois 群的表示记作 $\rho_\ell\colon G_F\to\GL(T_\ell E)$，称为 $E$ 的 $\ell$ 进表示。（有时这个名称也留给自同态代数的表示。）

$\ell$ 进表示的像是许多重要研究的主题。例如 Serre 证明：若 $F$ 是数域，且 $E$ 在 $F$ 上无复乘，则像在 $\GL(T_\ell E)\cong\GL(2,\mathbb{Z}_\ell)$ 中开，且对除有限多个 $\ell$ 以外的所有 $\ell$ 都是整个群。复乘情形较易理解，可追溯到 20 世纪上半叶。若 $E$ 在 $F$ 上有复乘，则 $\rho_\ell$ 的像交换。

更一般地，若把 $E$ 换成阿贝尔簇（自同态环可以更复杂），Tate 猜想、Faltings 证明：$G_F$ 的像的李代数是 $\operatorname{End}(E/\overline{F})$ 的整中心化子。这个深刻的定理是 Faltings 证明 Mordell 猜想的关键一步。另一方面，它是 Tate 关于代数圈的一大类猜想的一个例子。这个迷人的主题与我们的工作有些正交，就此打住。

\subsection{Weil 配对}\label{c5:sec:2.2}

Weil 发现了一个重要的、与 $G_F$ 及 $\operatorname{End}(E)$ 相容的、斜对称非退化配对
\begin{equation}\tag{2.3}\label{c5:eq:2.3}
\langle\cdot,\cdot\rangle\colon T_\ell E\times T_\ell E\longrightarrow T_\ell\mu,
\end{equation}
其中后者的意思是
\begin{equation}\tag{2.4}\label{c5:eq:2.4}
T_\ell\mu=\varprojlim\mu_{\ell^n},
\end{equation}
而 $\mu_r$ 是 $\overline{F}$ 中 $r$ 次单位根群。在复数上，若 $E=\mathbb{C}/\mathcal{L}$，其中
\begin{equation}\tag{2.5}\label{c5:eq:2.5}
\mathcal{L}=\mathbb{Z}+\mathbb{Z}\tau,
\end{equation}
则可用下述要求把它规范化：
\begin{equation}\tag{2.6}\label{c5:eq:2.6}
\langle \ell^{-n},\ell^{-n}\tau\rangle=e^{2\pi i/\ell^n}.
\end{equation}

\begin{annote}
Weil 配对 $e_N$ 的构造依赖 $g$ 的选取：把 $g$ 换成 $cg$，比值 $g(x)/g(x+b)$ 不变，故 $e_N$ 良定义。它满足 $e_N(a,b)=e_N(b,a)^{-1}$、$e_N(a,a)=1$、$e_N(a+b,c)=e_N(a,c)e_N(b,c)$，且非退化。复数上的规范化 (2.6) 把 $T_\ell E$ 的基取为 $(\ell^{-n},\ell^{-n}\tau)$：这两点在 $E[\ell^n]$ 中，配对值 $e^{2\pi i/\ell^n}$ 是标准的“$\ell^{-n}$ 次幂”约定。若把基换成 $(\ell^{-n}\tau,\ell^{-n})$，配对值变为其逆，这正是斜对称性的来源。
\par\vspace{0.15em}{\raggedleft 【来源：Silverman《The Arithmetic of Elliptic Curves》第 III 章（A）；Milne《Elliptic Curves》第 5 章（A）】\par}
\end{annote}

一般地，给定 $a\in E[N]$，存在函数 $f,g\in\overline{F}(E)$，其除子为
\begin{equation}\tag{2.7}\label{c5:eq:2.7}
\begin{aligned}
\operatorname{div}(f)&=N(a)-N(0),\\
\operatorname{div}(g)&=N^*((a)-(0))=\sum_{Na'=a}(a')-\sum_{Nu=0}(u).
\end{aligned}
\end{equation}
（这由 Abel--Jacobi 定理推出。）容易看出 $f\circ N=g^N$ 差一个常数，因为两者除子相同。若 $b\in E[N]$，定义
\begin{equation}\tag{2.8}\label{c5:eq:2.8}
e_N(a,b)=g(x)/g(x+b).
\end{equation}

这个函数是 $N$ 次单位根，这由 $g(x)^N=f(Nx)$ 推出。$T_\ell E$ 上的 Weil 配对由 $e_N$ 配对取 $N=\ell^n$ 的极限定义。然后可以证明它具有所有想要的性质。它是 Galois 等变的，且自同态 $\alpha$ 关于该配对的对偶是同源 $\check{\alpha}$。

Weil 配对对 $\ell$ 进表示存在性的直接推论是
\begin{equation}\tag{2.9}\label{c5:eq:2.9}
\det(\rho_\ell)=\chi_\ell,
\end{equation}
其中 $\chi_\ell$ 是分圆特征 $G_F\to\mathbb{Z}_\ell^\times$。若 $F=\mathbb{Q}$，由于 $\chi_\ell$ 把复共轭送到 $-1$，看出 $\rho_\ell$ 是奇的二维表示。

与 $\mathbb{Q}$ 张量损失很少：定义 $V_\ell=T_\ell E\otimes\mathbb{Q}\cong\mathbb{Q}_\ell^2$，并视
\begin{equation}\tag{2.10}\label{c5:eq:2.10}
\rho_\ell\colon G_F\longrightarrow\GL(V_\ell)\cong\GL(2,\mathbb{Q}_\ell)
\end{equation}
为表示，略去整格 $T_\ell E$ 的信息。当人们想把表示 $\rho_\ell$ 模 $\ell$ 约化时（这在 Wiles 的工作中至关重要），整格会再次起作用；但就构造 L 函数而言，它们无关紧要。

Weil 配对蕴含 $V_\ell$ 的反变对偶 $\check{V}_\ell$ 同构于 $V_\ell(-1)=V_\ell\otimes_{\mathbb{Z}_\ell}(T_\ell\mu)^\vee$。
\subsection{好归约与坏归约}\label{c5:sec:2.3}

从椭圆曲线 $E/\mathbb{Q}$ 开始。称 $E$ 在素数 $p$ 处有好归约，若存在局部环 $\mathbb{Z}_{(p)}$ 上的椭圆曲线 $\mathcal{E}$，其一般纤维为 $E$。一般环 $R$ 上椭圆曲线的定义与域上的类似：它是 $\Spec R$ 上固有、光滑的概形，几何纤维都是维数 $1$ 的，并带有一个截面 $s\colon\Spec R\to\mathcal{E}$。此时 $\mathcal{E}$ 唯一。特殊纤维 $\mathcal{E}\times_{\mathbb{Z}_{(p)}}\mathbb{F}_p$ 有时称为 $E$ 模 $p$ 的归约，记作 $\tilde{E}$。

若 $E$ 由上述 Weierstrass 平面模型给出，其中 $g_2,g_3$ 是 $p$ 进整数，且 $\Delta$ 不被 $p$ 整除，则显然 $E$ 在 $p$ 处有好归约，$\tilde{E}$ 由“把 Weierstrass 方程模 $p$ 读出来”得到。顺便说，这说明 $E$ 在几乎所有素数 $p$ 处有好归约。现在反过来也成立：若 $E$ 有好归约，则存在极小（广义）Weierstrass 方程，其 $\Delta$ 是 $p$ 进单位。然而，这关于 Weierstrass 模型数值的运气，不能推广到其他算术簇，甚至不能推广到更高亏格的曲线，所以好归约的概形论定义更好。

若 $E$ 没有好归约，称它有坏归约。人们仍可要求最佳可能的模型，它同样可用极小 Weierstrass 方程描述。这里的计算，尤其在特征 $2$ 与 $3$ 中，变得冗长。称 $E$ 有乘性归约（或半稳定归约），若存在固有平坦模型 $\mathcal{E}/\mathbb{Z}_{(p)}$，它不光滑，但其特殊纤维只含一个奇点，且该奇点是结点。否则称 $E$ 有加性归约，并可找到模型使其特殊纤维只有一个奇点，且该奇点是尖点。名称来自下述事实：在这些情形，特殊纤维上非奇点集构成一维的乘法群或加法群。

用极小 Weierstrass 方程来说，在特征不等于 $2$ 或 $3$ 时，乘性归约意味着三次方程的三个根中有两个模 $p$ 重合，加性归约意味着三个根模 $p$ 全都重合。

好归约或坏归约反映在由 $\rho_\ell$“切出”的（无限）扩张的分歧中。换言之，$\operatorname{Ker}(\rho_\ell)\subset G_{\mathbb{Q}}$ 是闭子群，其不动域是 $K_\ell=\mathbb{Q}(E[\ell^\infty])$，即把 $E$ 的全部 $\ell^n$ 扭点的坐标添加到 $\mathbb{Q}$ 得到的域。

\par\medskip\noindent\textbf{定理 2.3（Ogg--Néron--Shafarevich 判别法）}\quad 素数 $p\neq\ell$ 在 $\rho_\ell$ 中分歧（即在 $K_\ell$ 中分歧），当且仅当它是坏归约素数。

直到最近（过去 15 年内）才清楚如何通过观察表示 $\rho_p$ 来判断 $p$ 是好是坏；$\rho_p$ 在 $p$ 处总分歧（由 Weil 配对，$K_p$ 含 $p^n$ 次单位根），但这是另一段复杂得多的故事！

Ogg--Néron--Shafarevich 判别法的一个方向解释起来相对容易。若有好归约，$E[\ell^n]$ 中的点模 $p$ 上方素数单射（这涉及研究归约映射的核，它是 pro-$p$ 群，故不含 $\ell$ 扭）。但这意味着模 $p$ 上方素数的归约诱导同构
\begin{equation}\tag{2.11}\label{c5:eq:2.11}
T_\ell E\simeq T_\ell\tilde{E}.
\end{equation}
由于惯性群在剩余域上平凡作用，它在 $T_\ell E$ 上也平凡作用。另一方向更难，依赖 Néron 模型理论。事实上，有下述精确答案。回忆 $p\neq\ell$，$I_p$ 是 $K_\ell$ 中某个 $p$ 上方素数的惯性群。

\begin{itemize}
\item 若 $p$ 是好归约素数，则 $I_p=0$。
\item 若 $p$ 是乘性归约素数，则 $\dim(V_\ell)_{I_p}=1$。
\item 若 $p$ 是加性归约素数，则 $\dim(V_\ell)_{I_p}=0$。
\end{itemize}

\begin{annote}
表中三行常被误读成“惯性群越大，坏得越厉害”。实际含义是惯性不变子空间的维数：好归约时 $I_p$ 平凡作用，维数为 $0$；乘性归约时 $I_p$ 通过分圆特征作用在一维空间上，故不变子空间维数为 $1$；加性归约时 $I_p$ 可能非平凡，但 $(V_\ell)^{I_p}=0$。加性归约可以是“潜在好归约”，即 $I_p$ 在某个有限扩域上已平凡，此时 $V_\ell$ 作为 $I_p$ 模仍可能没有不变向量。判别法只涉及 $p\neq\ell$ 的情形；$p=\ell$ 时 $\rho_\ell$ 必在 $p$ 处分歧，需另作局部 Langlands 对应式的分析。
\par\vspace{0.15em}{\raggedleft 【来源：Silverman《The Arithmetic of Elliptic Curves》第 VII 章（A）；Serre--Tate《Good reduction of abelian varieties》（A）；nLab 条目“Néron model”（B，与 A 交叉核对）】\par}
\end{annote}

作为分类的最后一步，乘性情形还要分两种子情形。称 $E$ 有分裂乘性归约，若 $\tilde{E}$ 结点处的两条切线定义在 $\mathbb{F}_p$ 上。否则它们定义在二次扩张上，归约是非分裂的。Néron 模型理论让我们计算 Frobenius 类 $\sigma_p\in G_p/I_p$ 在 $(V_\ell)_{I_p}$ 上的作用。分裂情形它平凡作用，非分裂情形作用为 $-1$。

\section{Hasse--Weil L 函数}\label{c5:sec:3}

\subsection{定义}\label{c5:sec:3.1}

设 $E$ 是 $\mathbb{Q}$ 上的椭圆曲线。$E/\mathbb{Q}$ 的 L 函数由 $\ell$ 进表示 $\rho_\ell\colon G_{\mathbb{Q}}\to\GL(V_\ell E)$ 定义，配方与 Artin L 函数的定义相同。对每个素数 $p$ 选一个 $\ell\neq p$，并令
\begin{equation}\tag{3.1}\label{c5:eq:3.1}
L_p(E,s)=\det\bigl(1-p^{-s}\sigma_p\mid(V_\ell)_{I_p}\bigr)^{-1}.
\end{equation}
断言：这是 $p^{-s}$ 的整系数多项式的倒数，且与 $\ell$ 无关。

若 $p$ 是好归约素数，我们已看到 $V_\ell$ 无分歧，故可用 $V_\ell\tilde{E}$ 代替 $V_\ell E$。但前面注意到，在 $\mathbb{F}_p$ 上 $\sigma_p$ 的作用与 $\varphi_p\in\operatorname{End}(\tilde{E})$ 的相同。若 $\lambda$ 与 $\lambda'$ 是 $\varphi_p$ 特征多项式的两个根，习惯上写
\begin{equation}\tag{3.2}\label{c5:eq:3.2}
a_p=\operatorname{tr}(\varphi_p)=\lambda+\lambda',\qquad \lambda\lambda'=p.
\end{equation}
则 $a_p\in\mathbb{Z}$，且
\begin{equation}\tag{3.3}\label{c5:eq:3.3}
L_p(E,s)=(1-a_pp^{-s}+p^{1-2s})^{-1}.
\end{equation}
回忆 Hasse 的论文：
\begin{equation}\tag{3.4}\label{c5:eq:3.4}
|a_p|\le 2\sqrt{p}.
\end{equation}

若 $p$ 是乘性归约素数，我们已看到
\begin{equation}\tag{3.5}\label{c5:eq:3.5}
L_p(E,s)=(1-a_pp^{-s})^{-1},
\end{equation}
且 $a_p=1$（分裂情形）或 $-1$（非分裂情形）。

最后，若 $p$ 是加性归约素数，$a_p=0$，且
\begin{equation}\tag{3.6}\label{c5:eq:3.6}
L_p(E,s)=1.
\end{equation}

L 函数是 Euler 乘积
\begin{equation}\tag{3.7}\label{c5:eq:3.7}
L(E,s)=\prod_p L_p(E,s).
\end{equation}
它在 $\operatorname{Re}(s)>3/2$ 的紧集上一致收敛。这是 Hasse 建立的 $a_p$ 估计的推论。

\par\medskip\noindent\textbf{注记.}
\begin{enumerate}
\item[(a)] 这是附在 étale 上同调群上的 L 函数的第一个例子。一般地，好 Euler 因子与 $\ell$ 无关是 Deligne 证明 Weil 猜想的推论，但对坏 Euler 因子，这仍是未决问题。
\item[(b)] 好素数 $p$ 处的 Euler 因子与归约的 zeta 函数有关。后者定义为
\begin{equation}\tag{3.8}\label{c5:eq:3.8}
Z(\tilde{E},t)=\exp\Bigl(-\sum_{n=1}^{\infty}\frac{v_n}{n}t^n\Bigr),
\end{equation}
其中 $v_n=\#\tilde{E}(\mathbb{F}_{p^n})$。由前面得到的 $v_n$ 表达式，
\begin{equation}\tag{3.9}\label{c5:eq:3.9}
Z(\tilde{E},t)=\frac{(1-t)(1-pt)}{(1-\lambda t)(1-\lambda't)}
=\frac{(1-t)(1-pt)}{1-a_pt+pt^2},
\end{equation}
故
\begin{equation}\tag{3.10}\label{c5:eq:3.10}
Z(\tilde{E},p^{-s})=\frac{(1-p^{-s})(1-p^{1-s})}{L_p(E,s)}.
\end{equation}
\item[(c)] 当 $F=\mathbb{F}_q(X)$ 是有限域上光滑射影曲线的函数域，$E$ 是 $F$ 上椭圆曲线时，可用同样程序定义 $L(E,s)$，取遍整体域 $F$ 的全部素除子的 Euler 乘积。此时可给 $E$ 配一个 $X$ 上纤维化的曲面 $\mathcal{E}$，其一般纤维为 $E$。这个 $\mathcal{E}$ 可取为正则的（虽在 $X$ 上不光滑）。曲面的 zeta 函数与 $L(E,s)$ 之间于是有类似关系。这个关系，连同 Grothendieck 与 Deligne 关于有限域上代数簇 zeta 函数的工作，可用来证明 $L(E,s)$ 的解析延拓与函数方程；而当基域是数域 $F$ 时，这些仍是猜想（$F=\mathbb{Q}$ 除外！）。函数域上椭圆曲线的详细讨论见 Tate 的 Bourbaki 报告《On the conjecture of Birch and Swinnerton-Dyer and a geometric analogue》。
\item[(d)] 由于 $(V_\ell)_{I_p}$ 的对偶空间是 $\check{V}_\ell^{I_p}$，而 $\rho_\ell(\sigma_p)$ 的对偶是 $\check{\rho}_\ell(\sigma_p^{-1})$，又因对偶线性变换有相同的特征多项式，我们也可定义
\begin{equation}\tag{3.11}\label{c5:eq:3.11}
L_p(E,s)=\det\bigl(1-p^{-s}\sigma_p^{-1}\mid\check{V}_\ell^{I_p}\bigr)^{-1}.
\end{equation}
有些作者用这个。他们把 $\sigma_p^{-1}$ 称为几何 Frobenius。这里 $\check{V}_\ell^{I_p}=H^1(E_{\overline{\mathbb{Q}}},\mathbb{Q}_\ell)^{I_p}$（étale 上同调的惯性不变量）。
\end{enumerate}

\subsection{导子}\label{c5:sec:3.2}

与 Artin 导子的定义类似，可以给椭圆曲线 $E/\mathbb{Q}$ 附一个导子 $N$。它由 $\ell$ 进表示中的高阶分歧群决定。我们不给细节，只记录：

\begin{itemize}
\item 若 $E$ 在 $p$ 处好归约，则 $p\nmid N$。
\item 若 $E$ 在 $p$ 处乘性归约，则 $p\parallel N$。
\item 若 $E$ 在 $p$ 处加性归约且 $p\ge5$，则 $p^2\parallel N$。
\item 对 $p=2,3$，若 $E$ 加性归约，则 $p^2\mid N$，但 $N$ 中 $p$ 的幂可能更高。
\end{itemize}

\begin{annote}
导子指数 $f_p=v_p(N)$ 由惯性群在 $\ell$ 进表示上的分歧滤过（Swan 导子）控制。对 $p\ge5$ 的加性归约必有 $f_p=2$：此时潜在好归约，惯性像有限且温和，其作用分解为两个非平凡特征，各贡献 $1$。$p=2,3$ 时 $f_p$ 可能大于 $2$，需按 Ogg 公式逐例计算；书中“幂可能更高”即指此。乘性归约时 $f_p=1$，分裂与非分裂都一样。把这些指数与 (3.12) 的 $N^{s/2}$ 对照，可看出导子同时控制 Euler 因子的“坏程度”与函数方程中 $N$ 的幂次。
\par\vspace{0.15em}{\raggedleft 【来源：Silverman《The Arithmetic of Elliptic Curves》第 IV 章（A）；Ogg《Elliptic curves and wild ramification》（A）】\par}
\end{annote}

\subsection{完备 L 函数}\label{c5:sec:3.3}

域 $\mathbb{Q}$ 只有一个阿基米德位。定义
\begin{equation}\tag{3.12}\label{c5:eq:3.12}
\Lambda(E,s)=N^{s/2}(2\pi)^{-s}\Gamma(s)L(E,s).
\end{equation}

\par\medskip\noindent\textbf{猜想 3.1（Hasse--Weil）}\quad $\Lambda(E,s)$ 解析延拓到整个复平面，并满足函数方程
\begin{equation}\tag{3.13}\label{c5:eq:3.13}
\Lambda(E,s)=\pm\Lambda(E,2-s).
\end{equation}

\begin{annote}
伽马因子取 $\Gamma(s)$ 而非 $\Gamma(s/2)$，理由在 Mellin 变换：$\int_0^\infty f(iy)y^{s-1}\,dy=(2\pi)^{-s}\Gamma(s)L(f,s)$，对任意权都出现同一个 $\Gamma(s)$。权 $k$ 进入的是函数方程 $s\mapsto k-s$ 中的 $k$；对椭圆曲线 $k=2$，故 (3.13) 右端是 $\Lambda(E,2-s)$。若误以为权 $2$ 对应 $\Gamma(s/2)$，就会把函数方程的中心从 $s=1$ 挪到 $s=1/2$，整个归一化随之错位。因子 $N^{s/2}$ 同理是为了让 $N^{s/2}$ 与 $N^{(2-s)/2}$ 在 $s\mapsto2-s$ 下对齐。
\par\vspace{0.15em}{\raggedleft 【来源：Iwaniec--Kowalski《Analytic Number Theory》第 5 章（A）；Shimura《Introduction to the Arithmetic Theory of Automorphic Forms》（A）】\par}
\end{annote}

注意函数方程的符号应当决定 $\Lambda(E,s)$ 在 $s=1$ 处零点的阶的奇偶性。“特殊值” $L(E,1)$（更确切地说，$s=1$ 附近 Taylor 展开的首项）是算术中最有趣的特殊值之一。它应当记录 $E$ 的微妙算术不变量。只要说：猜想 $L(E,s)$ 在 $s=1$ 处零点的阶等于 $E(\mathbb{Q})$ 的 Mordell--Weil 秩（Birch 与 Swinnerton-Dyer 猜想）就够了！

$\mathbb{Q}$ 上椭圆曲线的 Hasse--Weil 猜想现在是定理，归功于上面提到的 Wiles 等人的工作。但证明是间接的。它基于模形式的 L 函数满足这些性质，以及 Shimura--Taniyama--Weil 猜想的证明；该猜想断言对适当的模形式 $f$ 有 $L(E,s)=L(f,s)$。

在一般数域上，Hasse--Weil 猜想仍远未解决。当然，$L(E,s)$ 的定义与函数方程的形式都要修改，但 $\mathbb{Q}$ 与一般数域之间的差距不是技术性的，而与 Wiles 纲领目前的局限有关。

用任意光滑射影曲线 $C$（定义在 $\mathbb{Q}$ 上）的雅可比簇及其 $\ell$ 进表示（此时维数为 $2g$，$g=\mathrm{genus}(C)$），可对每条曲线附上 L 函数 $L(C,s)$，方式与上面类似。$L(C,s)$ 的好 Euler 因子次数为 $2g$（作为 $p^{-s}$ 的多项式）。一般地，若 $g>1$，即使基域是 $\mathbb{Q}$，$L(C,s)$ 的解析延拓与函数方程仍未知。

\subsection{具有复乘的椭圆曲线}\label{c5:sec:3.4}

还有一类情形，Hasse--Weil 猜想在任意数域上早已知道，即 $E$ 具有复乘的情形。为简单起见，仍设 $E$ 定义在 $\mathbb{Q}$ 上，且在 $\overline{\mathbb{Q}}$ 上具有二次虚域 $K$ 的整数环 $\mathcal{O}_K$ 的复乘。此时自同态已经定义在 $K$ 上。我们规范化同构
\begin{equation}\tag{3.14}\label{c5:eq:3.14}
\mathcal{O}_K\simeq\operatorname{End}(E/K),
\end{equation}
使得自同态 $\alpha$ 在 $0$ 处切空间上的作用是乘以 $\alpha$（而非 $\bar\alpha$）。

复乘理论表明，此时把 $E$ 的扭点添加到 $K$ 所生成的扩张在 $K$ 上阿贝尔。给 $E$ 配一个 $K$ 的 Hecke 特征（Größencharakter）$\psi$。它是从 $I(\mathfrak{f})$（与 $\mathfrak{f}$ 互素的分数理想群）到 $K^\times$ 的特征，在此情形，对 $\alpha\equiv1\bmod\mathfrak{f}$ 满足
\begin{equation}\tag{3.15}\label{c5:eq:3.15}
\psi((\alpha))=\alpha
\end{equation}

其定义性质如下：若 $\mathfrak{p}$ 是 $K$ 的不整除 $\mathfrak{f}$ 的素理想，$\ell$ 与 $\mathfrak{p}$ 互素，则 $K_\ell=K(E[\ell^\infty])$ 在 $\mathfrak{p}$ 处无分歧，且 $\mathfrak{p}$ 的 Frobenius 自同构 $\sigma_{\mathfrak{p}}$ 通过 $\psi(\mathfrak{p})$ 作用：
\begin{equation}\tag{3.16}\label{c5:eq:3.16}
\rho_\ell(\sigma_{\mathfrak{p}})(x)=\psi(\mathfrak{p})(x)\qquad\text{对 }x\in T_\ell E.
\end{equation}

$\psi$ 的导子 $\mathfrak{f}$ 与 $E/\mathbb{Q}$ 的导子 $N$ 由下述关系联系：
\begin{equation}\tag{3.17}\label{c5:eq:3.17}
N=d_KN\mathfrak{f}.
\end{equation}

$\psi$ 的 Hecke L 函数定义为对 $K$ 的全部理想求和：
\begin{equation}\tag{3.18}\label{c5:eq:3.18}
L(\psi,s)=\sum_{\mathfrak{a}}\psi(\mathfrak{a})\mathrm{N}\mathfrak{a}^{-s}.
\end{equation}

比较 Euler 因子得到等式
\begin{equation}\tag{3.19}\label{c5:eq:3.19}
L(E/\mathbb{Q},s)=L(\psi,s).
\end{equation}

\begin{annote}
(3.14) 的标准化“乘 $\alpha$ 而非 $\bar\alpha$”是关键：它把 $\mathcal{O}_K$ 与 $\operatorname{End}(E/K)$ 的生成元对齐，使 (3.15) 的 $\psi((\alpha))=\alpha$ 成立。若换成 $\bar\alpha$，得到的 $\bar\psi$ 只与 $\psi$ 相差复共轭，而 $L(\bar\psi,s)=\overline{L(\psi,s)}$，解析性质相同。关系 (3.17) 中的 $d_K$ 是 $K$ 的判别式：$L(\psi,s)$ 作为 $\mathbb{Q}$ 上的 L 函数，其导子必须是 $\mathbb{Q}$ 的理想，$d_K$ 正是把 $\mathfrak{f}$ 的范数“下降”到 $\mathbb{Q}$ 的修正项。复乘情形之所以能绕过 Wiles，正因为 Hecke 早已对任意数域证明了 $L(\psi,s)$ 的解析延拓与函数方程。
\par\vspace{0.15em}{\raggedleft 【来源：Shimura《Introduction to the Arithmetic Theory of Automorphic Forms》第 7 章（A）；Silverman《Advanced Topics in the Arithmetic of Elliptic Curves》第 II 章（A）】\par}
\end{annote}

Hecke 已建立（对任意 Hecke 特征与任意基域）$L(\psi,s)$ 的解析延拓与函数方程。这将在 Kudla 关于 Tate 论文的一章中解释。要点是：Hecke 特征，如同有限阶广义理想类群的特征，可解释为伊代尔类群特征。对 $E$ 而言，这就归结为 $L(E,s)$ 的 Hasse--Weil 猜想。
\section{模曲线及其雅可比簇}\label{c5:sec:4}

本节假定读者熟悉模形式、其 L 函数与 Hecke 算子的解析理论，如 Kowalski 所述。

\subsection{模曲线}\label{c5:sec:4.1}

复模曲线 $X_0(N)$ 定义为开 Riemann 曲面
\begin{equation}\tag{4.1}\label{c5:eq:4.1}
Y_0(N)=\Gamma_0(N)\backslash\mathfrak{H}
\end{equation}
的紧化，其中 $\mathfrak{H}$ 是上半平面，且
\begin{equation}\tag{4.2}\label{c5:eq:4.2}
\Gamma_0(N)=\Bigl\{\begin{pmatrix}a&b\\ c&d\end{pmatrix}\in\mathrm{SL}(2,\mathbb{Z});\ c\equiv0\bmod N\Bigr\}.
\end{equation}

$X_0(N)-Y_0(N)$ 的点称为尖点，在模形式理论中起特殊作用。

考虑把 $\tau\in\mathfrak{H}$ 送到偶
\begin{equation}\tag{4.3}\label{c5:eq:4.3}
(E_\tau,H_\tau)
\end{equation}
的映射，其中 $E_\tau=\mathbb{C}/\mathcal{L}_\tau$，$\mathcal{L}_\tau=\mathbb{Z}\tau+\mathbb{Z}$，$H_\tau=\langle\frac1N\bmod\mathcal{L}_\tau\rangle$。这个偶由一条椭圆曲线及其 $N$ 阶循环子群组成。两个这样的偶的同构是椭圆曲线的同构，并把循环子群彼此对应。于是容易验证：$\tau$ 与 $\tau'$ 对应的偶同构，当且仅当 $\tau$ 与 $\tau'$ 在 $\Gamma_0(N)$ 的同一轨道中。因此 $Y_0(N)$ 是上述偶 $(E,H)$ 的模空间。适当定义“广义（或退化）椭圆曲线”，也可给尖点一个模解释。

\begin{annote}
$H_\tau=\langle\frac1N\bmod\mathcal{L}_\tau\rangle$ 依赖基的选取：$\mathcal{L}_\tau=\mathbb{Z}\tau+\mathbb{Z}$ 的基取作 $(\tau,1)$，故 $\frac1N$ 指 $\frac1N\cdot1$。$\Gamma_0(N)$ 恰是使偶 $(E_\tau,\langle\frac1N\rangle)$ 与 $(E_{\gamma\tau},\langle\frac1N\rangle)$ 同构的 $\gamma$ 全体；若改用 $\Gamma^0(N)$（要求 $b\equiv0$），对应的循环子群要换成 $\langle\frac{\tau}{N}\rangle$。两套水平结构给出同构的模曲线，但尖点的标号不同。尖点对应“广义椭圆曲线”，即 Néron 多边形退化的情形，须按 Deligne--Rapoport 的意义理解。
\par\vspace{0.15em}{\raggedleft 【来源：Diamond--Shurman《A First Course in Modular Forms》第 8 章（A）；Deligne--Rapoport《Les schémas de modules de courbes elliptiques》（A）】\par}
\end{annote}

代数几何提供了研究这类模问题的工具，可在任意环上进行，并用代数簇或概形表示它们。人们正是这样得到 $X_0(N)$ 的典范模型：先是在 $\mathbb{Q}$ 上（Kronecker、Shimura），然后在 $\mathbb{Z}[\frac1N]$ 上（Igusa），最后在 $\mathbb{Z}$ 上（Deligne--Rapoport）。模曲线 $X_0(N)$，以及由 $\mathrm{SL}(2,\mathbb{Z})$ 的其他同余子群构造的类似模曲线，是迷人的对象。由于模解释，它们带有大量对称性——更确切地说，带有一个大的对应环（Hecke 代数）——这使其几何更丰富。由于它们至少在 $\mathbb{C}$ 上以群论方式定义——回忆 $\mathfrak{H}=\mathrm{SL}(2,\mathbb{R})/\mathrm{SO}(2)$——其 L 函数与自守 L 函数重合，而后者可证解析延拓与函数方程。

我们关心 $X_0(N)$ 与导子为 $N$ 的椭圆曲线之间的关系，这初看相当出人意料。开头要澄清一个易混点：这个关系（即本章的主题）与 $X_0(N)$ 自身是带水平结构的椭圆曲线模空间这一事实毫无关系。于是椭圆曲线以两种完全无关的方式介入！

对模形式的研究而言，模曲线重要，因为它们是模形式所生活的空间。更确切地说，模形式的正确几何解释——它有许多算术推论（如 Fourier 系数增长的 Ramanujan 估计）——是模曲线上某些线丛的截面。本章只关心权 $2$ 的模形式。若
\begin{equation}\tag{4.4}\label{c5:eq:4.4}
f\in S_2(\Gamma_0(N))
\end{equation}
是 $\Gamma_0(N)$ 上的尖点形式，则
\begin{equation}\tag{4.5}\label{c5:eq:4.5}
\omega_f=2\pi if(z)\,dz
\end{equation}
成为 $\mathfrak{H}$ 上 $\Gamma_0(N)$ 不变微分形式，即 $X_0(N)$ 上的全纯微分形式。（尖点条件恰好保证 $\omega_f$ 在尖点处也全纯。）众所周知，Riemann 曲面 $X$ 上全纯微分形式空间 $H^0(X,\Omega)$ 的维数为 $g=\mathrm{genus}(X)$，并可等同于 $X$ 的雅可比簇 $\operatorname{Jac}(X)$ 在 $0$ 处的余切空间；$\operatorname{Jac}(X)$ 是维数 $g$ 的复环面。因此我们转而研究 $J_0(N)$，即 $X_0(N)$ 的雅可比簇，及其自同态环。

\subsection{模雅可比簇}\label{c5:sec:4.2}

刚说过 $J_0(N)$ 是维数 $g$ 的复环面。此外它具有光滑射影群簇——阿贝尔簇——的结构，且定义在 $\mathbb{Q}$ 上。因此它有 Tate 模，类似于椭圆曲线的 Tate 模：
\begin{equation}\tag{4.6}\label{c5:eq:4.6}
T_\ell J_0(N)=\varprojlim J_0(N)[\ell^n]\cong\mathbb{Z}_\ell^{2g},
\end{equation}
$G_{\mathbb{Q}}$ 与 $\operatorname{End}(J_0(N))$ 在其上交换作用。对某些小的 $N$，例如 $N=11$，$g=1$，$X_0(N)=J_0(N)$ 本身就是椭圆曲线。一般地，这是我们对椭圆曲线详细讨论过的高维类比。

由于 $X_0(N)$ 在 $\mathbb{Z}[\frac1N]$ 上有好归约，其雅可比簇也是。由 Ogg--Néron--Shafarevich 判别法，$\ell$ 进表示
\begin{equation}\tag{4.7}\label{c5:eq:4.7}
\rho_\ell\colon G_{\mathbb{Q}}\longrightarrow\GL(V_\ell J_0(N))
\end{equation}
在每个不整除 $N\ell$ 的素数处无分歧。

\subsection{Hecke 代数}\label{c5:sec:4.3}

两条光滑射影曲线 $X$ 与 $X'$ 之间的对应是曲面 $X\times X'$ 上的一条（可以可约，但要完全）曲线 $\Gamma$，它不含竖直或水平分支。应把 $\Gamma$ 想成从 $X$ 到 $X'$（或反向）的多值函数的图像。对应在形式加法下自然构成阿贝尔群；若 $X=X'$，以复合为乘法，它们构成环（把对应想成多值函数并用通常的函数复合，复合最易理解）。$X$ 与 $X'$ 之间任一对应 $\Gamma$ 诱导阿贝尔簇的同态
\begin{equation}\tag{4.8}\label{c5:eq:4.8}
\gamma\colon\operatorname{Jac}(X)\longrightarrow\operatorname{Jac}(X').
\end{equation}

例如，把雅可比簇等同于 Picard 簇，设 $p$ 与 $p'$ 是 $\Gamma$ 到 $X$ 与 $X'$ 的两个投影，取 $X$ 上次数 $0$ 的除子 $D$，记 $[D]$ 为其模主除子的类，定义
\begin{equation}\tag{4.9}\label{c5:eq:4.9}
\gamma([D])=[p'_*p^*D].
\end{equation}
这是良定义的。此外 $\Gamma\mapsto\gamma$ 可加；当 $X=X'$ 时，它是 $X$ 的对应环到 $\operatorname{End}(\operatorname{Jac}(X))$ 的环同态。

由于模解释，模曲线 $X_0(N)$ 有许多有趣的对应。对每个不整除 $N$ 的素数 $p$，取点 $x\in X_0(N)$，用 $(E,H)$ 表示它，设 $C_0,\dots,C_p$ 是 $E$ 的 $p+1$ 个 $p$ 阶循环子群。设 $T_p$ 是把 $x$ 送到形式和 $x_0+\cdots+x_p$ 的对应，其中 $x_i$ 由 $(E/C_i,H+C_i\bmod C_i)$ 表示。稍加小心，对 $p\mid N$ 也可定义类似对应。这些 $T_p$ 称为 Hecke 对应。它们在 $J_0(N)$ 上，或经函子性在其余切空间（即权 $2$ 尖点形式空间）上诱导的自同态，称为 Hecke 算子。它们对尖点形式 $q$ 展开的作用由 Kowalski 一章中记录的公式给出。

它们生成的环 $\mathbb{Z}[\dots,T_p,\dots]$ 称为 Hecke 代数。它交换。对一般的 $N$ 它可能不半单，这一事实导致 Atkin--Lehner 的新旧形式理论。为避免这一技术复杂性，从现在起（至少涉及 Hecke 代数时）假定
\begin{equation}\tag{4.10}\label{c5:eq:4.10}
N\ \text{是素数}.
\end{equation}
这意味着（有理）Hecke 代数是半单的，是域的乘积，因此 $S_2(\Gamma_0(N))$ 有一组同时为所有 Hecke 算子本征形式的本征形式基。

\section{Eichler--Shimura 同余关系}\label{c5:sec:5}

\subsection{同余关系的证明}\label{c5:sec:5.1}

若 $p$ 是 $X=X_0(N)$ 的好素数（不整除 $N$ 的素数），则 $X$ 延拓为 $\mathbb{Z}_{(p)}$ 上的光滑射影曲线 $\tilde{X}$，它在 $\mathbb{Z}_{(p)}$ 代数上分类同样的偶 $(E,H)$。特别地可考虑它模 $p$ 的归约，仍记作 $\tilde{X}$。对应 $T_\ell$ 作为 $X\times X$ 上的曲线，也经受类似的归约 $\tilde{T}_\ell$。我们希望用特征 $p$ 中模问题上的对应来理解它。对 $\ell\neq p$ 没有困难。但对 $T_p$ 有问题：特征 $p$ 中的椭圆曲线没有 $p+1$ 个 $p$ 阶循环子群。若 $E$ 是常的，只有一个；若 $E$ 超奇异，一个也没有！

为克服困难，应当谈论 $E$ 中阶为 $p$ 的有限群概形。特征 $p$ 中总有一个特殊的，即次数 $p$ 的 Frobenius 同源的核：
\begin{equation}\tag{5.1}\label{c5:eq:5.1}
\varphi_E\colon E\longrightarrow E^{(p)},
\end{equation}
其中 $E^{(p)}$ 是把定义 $E$ 的方程系数提升到 $p$ 次幂得到的曲线。（若 $E$ 定义在 $\mathbb{F}_q$ 上，$q=p^f$，则 $\varphi_E$ 的第 $f$ 次迭代就是前面考虑过的 $E$ 的自同态。）若 $E$ 是常的（这是通有情形），对偶同源
\begin{equation}\tag{5.2}\label{c5:eq:5.2}
E^{(p^{-1})}\longleftarrow E:\tilde{\varphi}_{E^{(p^{-1})}}
\end{equation}
可分，其核是 $E$ 的“真正的”（确切说，既约的）$p$ 阶循环子群。若 $E$ 超奇异，$E^{(p^{-1})}=E^{(p)}$，且 $\tilde{\varphi}_{E^{(p^{-1})}}$ 在差一个曲线自同构的意义下等于 $\varphi_E$，故两者核相同。无论如何，这些是特征 $p$ 中仅有的（$p$ 阶）子群概形。

\begin{annote}
记号 $E^{(p^{-1})}$ 需要解释：若 $E$ 由系数在 $\mathbb{F}_q$ 中的方程给出，则 $E^{(p)}$ 是把系数提升到 $p$ 次幂得到的曲线，而 $E^{(p^{-1})}$ 是满足 $(E^{(p^{-1})})^{(p)}=E$ 的那条曲线，即把系数开 $p$ 次方。$\varphi_E\colon E\to E^{(p)}$ 的对偶同源方向相反，记作 $\tilde{\varphi}_{E^{(p^{-1})}}\colon E\to E^{(p^{-1})}$，它正是曲线层面的 Verschiebung $V$。常曲线时它可分、核为 $p$ 阶循环子群；超奇异时 $E^{(p^{-1})}=E^{(p)}$，它与 $\varphi_E$ 只差自同构，核相同。(5.2) 的箭头方向正是这一点的体现。
\par\vspace{0.15em}{\raggedleft 【来源：Silverman《The Arithmetic of Elliptic Curves》第 V 章（A）；Mumford《Abelian Varieties》第 12 节（A）】\par}
\end{annote}

若 $E$ 是特征 $0$ 的椭圆曲线（代表 $X$ 上的点 $x$），且在 $p$ 处好归约（意思是 $x$ 不归约到 $\tilde{X}$ 的尖点），则 $p+1$ 个 $p$ 阶子群中有一个，比如 $C_0$，在 $p$ 进拓扑下靠近原点，并归约到 $\operatorname{Ker}(\varphi_{\tilde{E}})$。其余 $p$ 个子群都归约到 $\operatorname{Ker}\tilde{\varphi}_{\tilde{E}^{(p^{-1})}}$。于是除以 $C_i$ 后，我们一次得到 $\tilde{E}^{(p)}$，$p$ 次得到 $\tilde{E}^{(p^{-1})}$。对辅助的 $N$ 级结构同样如此。习惯上把 $\tilde{X}$ 上把 $x=[(E,H)]$（特征 $p$ 中）送到 $x^{(p)}$ 的对应称为 Frobenius 对应，记作 $F$。对偶对应（交换源与靶的角色得到）称为 Verschiebung（转移），记作 $V$。用前面的记号，它们是由
\begin{equation}\tag{5.3}\label{c5:eq:5.3}
\Gamma=\bigl\{(x,x^{(p)})\in\tilde{X}\times\tilde{X}\bigr\}
\end{equation}
得到的两个对应。

注意它们特殊于特征 $p$，不能提升到特征 $0$。也注意第一投影 $\Gamma\to\tilde{X}$ 的次数为 $1$（是同构），但第二投影不可分、次数为 $p$。故 $F$ 的次数为 $1$，$V$ 纯不可分、次数为 $p$。在点上的作用为
\begin{equation}\tag{5.4}\label{c5:eq:5.4}
F:x\longmapsto x^{(p)},\qquad V:x\longmapsto p\,x^{(p^{-1})}.
\end{equation}

我们已勾勒了下述定理的证明。

\par\medskip\noindent\textbf{定理 5.1（Eichler--Shimura 同余关系）}\quad 对应 $T_p$ 模 $p$ 的归约是
\begin{equation}\tag{5.5}\label{c5:eq:5.5}
\tilde{T}_p=F+V.
\end{equation}

虽然 $F$ 与 $V$ 各自不能提升到特征 $0$，它们的和却能！这个重要定理的种子已见于 Kronecker 关于模方程的工作。

\begin{annote}
(5.5) 的分量 $F$ 与 $V$ 各自都不是特征 $0$ 中对应的归约，但它们的和是。原因是 $T_p$ 的“$p+1$ 个循环子群”在特征 $p$ 中退化：一个子群变成 $\ker\varphi_{\tilde{E}}$（即 $F$），其余 $p$ 个挤成 $\ker\tilde{\varphi}_{\tilde{E}^{(p^{-1})}}$（即 $V$），重数恰好 $1$ 与 $p$。在 $T_\ell J_0(N)$ 上，$F$ 的特征值是 $\lambda_i,\bar\lambda_i$，$V=pF^{-1}$ 的特征值是 $\bar\lambda_i,\lambda_i$，故 $F+V$ 的特征值是 $\lambda_i+\bar\lambda_i=a_{p,i}$，每个重数 $2$。Weil 界 $|\lambda_i|=\sqrt{p}$ 保证这些值落在 Hasse 区间内。
\par\vspace{0.15em}{\raggedleft 【来源：Shimura《Introduction to the Arithmetic Theory of Automorphic Forms》第 7 章（A）；Deligne《Formes modulaires et représentations $\ell$-adiques》（A）】\par}
\end{annote}

\subsection{对模曲线 L 函数的推论}\label{c5:sec:5.2}

正如我们在椭圆曲线情形观察到的，要研究 $G_{\mathbb{Q}}$ 中 Frobenius 自同构 $\sigma_p$ 在 $J_0(N)$ 的 $\ell$ 进表示中的作用（$p$ 好），可模 $p$ 约化，于是算术定义的 $\sigma_p$ 变成几何定义的 $F$。Weil 1948 年关于有限域上曲线 Riemann 假设的工作（推广 Hasse 关于椭圆曲线的论文！）告诉我们：$F$ 在 $T_\ell J_0(N)$ 上作用的特征值是绝对值为 $\sqrt{p}$（在每个复嵌入下）的代数整数 $\lambda_1,\dots,\lambda_g,\bar\lambda_1,\dots,\bar\lambda_g$。$V$（在同一基下）的特征值是 $\bar\lambda_1,\dots,\bar\lambda_g,\lambda_1,\dots,\lambda_g$。于是 $T_p$ 的特征值是 $a_{p,1},\dots,a_{p,g}$，每个重数为 $2$，其中 $a_{p,i}=\lambda_i+\bar\lambda_i$。

$T_p$ 在 $T_\ell J_0(N)$ 上的作用同构于它在雅可比簇奇异同调上的作用，后者是它在余切空间及其复共轭上的表示之和。但余切空间（作为 Hecke 模）等同于 $S_2(\Gamma_0(N))$，且 $T_p$ 在尖点形式上的特征值是实数。因此我们断定 $a_{p,1},\dots,a_{p,g}$ 正是 $T_p$ 在 $S_2(\Gamma_0(N))$ 上作用的特征值。

$X$ 在 $\mathbb{Q}$ 上的 L 函数在 $p$ 处的 Euler 因子是
\begin{equation}\tag{5.6}\label{c5:eq:5.6}
\begin{gathered}
\prod_{i=1}^{g}(1-\lambda_ip^{-s})(1-\bar\lambda_ip^{-s})\\
=\prod_{i=1}^{g}(1-a_{p,i}p^{-s}+p^{1-2s}).
\end{gathered}
\end{equation}

\par\medskip\noindent\textbf{推论 5.2（Shimura，Carayol）}\quad 设 $f_1,\dots,f_g$ 是 $S_2(\Gamma_0(N))$ 的一组基，由 Hecke 算子的规范本征函数组成（回忆我们现在假定 $N$ 为素数，故所有 $f_i$ 都是新形式）。设
\begin{equation}\tag{5.7}\label{c5:eq:5.7}
f_i=\sum_{n=1}^{\infty}a_{n,i}q^n
\end{equation}
是它们的 $q$ 展开，故 $T_pf_i=a_{p,i}f_i$。设
\begin{equation}\tag{5.8}\label{c5:eq:5.8}
L(f_i,s)=\sum_n a_{n,i}n^{-s}=\prod_p(1-a_{p,i}p^{-s}+p^{1-2s})^{-1}
\end{equation}
是相应的 L 函数。则
\begin{equation}\tag{5.9}\label{c5:eq:5.9}
L(X_0(N),s)=\prod_i L(f_i,s).
\end{equation}

好 Euler 因子（$p$ 不整除 $N$）相符，正如我们刚验证的。坏 Euler 因子也相符，但这是 Carayol 后来才证明的，且更难。注意要建立解析延拓，只需知道 L 函数到有限多个 Euler 因子为止。只有函数方程的精确形式才需要知道全部 Euler 因子。

推论的一个重要结果是：模曲线 $X_0(N)$ 的 Hasse--Weil 猜想成立。这当然依赖 Kowalski 一章中讨论的 $L(f_i,s)$ 的解析延拓与函数方程。现在我们准备考虑 $\mathbb{Q}$ 上的椭圆曲线。
\section{椭圆曲线的一致化与 Shimura--Taniyama--Weil 猜想}\label{c5:sec:6}

\subsection{作为 $J_0(N)$ 商的椭圆曲线}\label{c5:sec:6.1}

为简化叙述，保留 $N$ 为素数的假定。空间 $S_2(\Gamma_0(N))$ 有一个很好的解释：从 Hecke 代数到 $\mathbb{C}$ 的线性映射空间。本征形式正是那些同时也是代数同态的线性映射。设 $f=\sum a_nq^n$ 是规范化的尖点形式，是所有 Hecke 算子的本征形式，并设其 Fourier 系数有理：$a_n\in\mathbb{Z}$。把 $f$ 解释为代数同态，看出有理 Hecke 代数
\begin{equation}\tag{6.1}\label{c5:eq:6.1}
\mathcal{H}=\mathbb{Q}[\dots,T_p,\dots]\subset\operatorname{End}(J_0(N))\otimes\mathbb{Q}
\end{equation}
有一个同构于 $\mathbb{Q}$ 的商，其中 $T_p$ 映到 $a_p$。相应的幂等元 $e_f$ 是有理自同态代数中的投影；它的某个整数倍是真正的自同态，把 $J_0(N)$ 投影到一条椭圆曲线商上（维数可如常从余切空间确定）。称之为 $E_f$。它定义在 $\mathbb{Q}$ 上，因为 $J_0(N)$ 与投影 $e_f$ 都定义在 $\mathbb{Q}$ 上。

上面关于整个模雅可比簇的论证的变体，现在蕴含 $a_p$（对好的 $p$）是 Frobenius 在 $E_f$ 上的迹。如常，坏 Euler 因子也相符，尽管这更难。我们得到下述定理。

\par\medskip\noindent\textbf{定理 6.1}\quad $L(E_f,s)=L(f,s)$ 满足 Hasse--Weil 猜想。

Ogg 证明了 $E_f$ 的导子是 $N$。（算术导子等于解析导子。）

\subsection{Shimura--Taniyama--Weil 猜想}\label{c5:sec:6.2}

50 年代后期的某个时候，下述猜想从日本传到 Princeton，引起 A.~Weil 的注意：

\par\medskip\noindent\textbf{猜想 6.2}\quad $\mathbb{Q}$ 上每条椭圆曲线 $E$ 都同源于某个 $E_f$。更确切地说，若 $N$ 是 $E$ 的导子，则存在水平 $N$、权 $2$ 的新形式，使得 $E$ 同源于 $E_f$。

起初这个猜想证据很少，没人敢把它写下来。但 1967 年 Weil 证明了他的逆定理（在我们前面一章讨论过）。它有以下推论。

\par\medskip\noindent\textbf{定理 6.3（Weil）}\quad 设 $L(E,s)=\sum a_nn^{-s}$，连同所有扭
\begin{equation}\tag{6.2}\label{c5:eq:6.2}
L(E,\chi,s)=\sum a_n\chi(n)n^{-s},
\end{equation}
对每个导子与某整数 $M$ 互素的 Dirichlet 特征 $\chi$，都满足解析延拓与函数方程（类似为 $L(E,s)$ 预言的那个）。则 $E=E_f$ 是模的。

与 Artin L 函数的情形一样，$E$ 的模性（或所附 $\ell$ 进表示的模性，或 L 函数的模性，三者等价）成为所有 $L(E,\chi,s)$ 的 Hasse--Weil 猜想的充要条件。

1995 年 Taylor 与 Wiles 对无平方因子的 $N$ 等情形证明了猜想，但并非全部情形。Diamond、Taylor、Conrad、Breuil 的进一步工作消除了对 $N$ 的限制，所以今天猜想对所有 $N$ 都已证明。

有趣的是，这个猜想可纯粹表述为复函数论中的猜想：$E$ 同源于 $J_0(N)$ 的一个因子，当且仅当存在从 $X_0(N)$ 到 $E$ 的非恒定 Riemann 曲面映射。

\begin{annote}
“同源”不能减弱为“同构”：给定 $E$，与之同源的 $E_f$ 一般不唯一，但同源类唯一。等价的几何说法是存在非恒定映射 $X_0(N)\to E$，这样的映射经 Abel--Jacobi 嵌入对偶化为 $J_0(N)\to E$ 的商，其核是阿贝尔簇；书中此说法里的 $N$ 就是 $E$ 的导子。导子条件不能省：$E_f$ 的导子由 Ogg 定理等于 $N$。历史地，Taylor--Wiles 1995 年处理半稳定（含无平方因子导子）情形，Breuil--Conrad--Diamond--Taylor 2001 年完成一般情形；书中“today”指 2004 年成书时。
\par\vspace{0.15em}{\raggedleft 【来源：Darmon--Diamond--Taylor《Fermat's last theorem》综述（A）；Cornell--Silverman--Stevens 编《Modular Forms and Fermat's Last Theorem》（A）】\par}
\end{annote}

我们不试图阐述 Taylor 与 Wiles 的工作。已有极好的综述文章，但它们所需的背景都超出本书想要假定的范围。

\section*{参考文献}\label{c5:sec:ref}

关于这两章的主题，已有许多优秀教科书。我们列出少数几种。关于 Fermat 大定理的参考文献也含有关于前述主题的极好综述章节。

\par\medskip\noindent\textbf{关于椭圆曲线：}
\begin{enumerate}[label={[\arabic*]},leftmargin=2.2em,itemsep=0.15em]
\item A.~Robert, \emph{Elliptic Curves}, Springer-Verlag, New York, 1973.
\item S.~Lang, \emph{Elliptic Functions}, Addison--Wesley and Springer-Verlag, New York and Reading, MA, 1987.
\item J.~Silverman, \emph{The Arithmetic of Elliptic Curves}, Springer-Verlag, New York, 1986.
\item J.~Silverman, \emph{Advanced Topics in the Arithmetic of Elliptic Curves}, Springer-Verlag, New York, 1994.
\end{enumerate}

\par\medskip\noindent\textbf{关于 $\ell$ 进表示：}
\begin{enumerate}[label={[\arabic*]},leftmargin=2.2em,itemsep=0.15em,start=5]
\item J.-P.~Serre, \emph{Abelian $\ell$-Adic Representations and Elliptic Curves}, Benjamin, New York, 1968.
\end{enumerate}

\par\medskip\noindent\textbf{关于模形式、模曲线及其与椭圆曲线的关系：}
\begin{enumerate}[label={[\arabic*]},leftmargin=2.2em,itemsep=0.15em,start=6]
\item G.~Shimura, \emph{Introduction to the Arithmetic Theory of Automorphic Forms}, Princeton University Press, Princeton, NJ, 1971.
\item S.~Lang, \emph{Introduction to Modular Forms}, Springer-Verlag, New York, 1976.
\item A.~Knapp, \emph{Elliptic Curves}, Princeton University Press, Princeton, NJ, 1992.
\end{enumerate}

\par\medskip\noindent\textbf{关于 Wiles 对 Taniyama--Weil 猜想的工作：}
\begin{enumerate}[label={[\arabic*]},leftmargin=2.2em,itemsep=0.15em,start=9]
\item H.~Darmon, F.~Diamond, and R.~Taylor, Fermat's last theorem (a survey), \emph{Curr. Dev. Math.}, \textbf{1}(1995), 1--157.
\item G.~Cornell, J.~Silverman, and G.~Stevens, eds., \emph{Modular Forms and Fermat's Last Theorem}, Springer-Verlag, New York, 1997.
\item K.~Murty, ed., \emph{Seminar on Fermat's Last Theorem}, CMS Conference Proceedings 17, American Mathematical Society, Providence, 1995.
\end{enumerate}
